\section{Introduction (Reuben)}
\textbf{First paragraph:}
\begin{itemize}
    \item Context: Diffusion MRI provides quantitative microstructural biomarkers
    \item Increasing use of ML for prediction (age, diagnosis, cognition, etc.)
    \item Lack of standardized benchmark
\end{itemize}

\noindent \textbf{Second paragraph: literature review:}
\begin{itemize}
    \item  Literature review of techniques: Variability of the proposed approach (ML-based, deep-learning based etc)
    \item Problems: High variability in reported performance across studies + Limited reproducibility 
\end{itemize}

\noindent \textbf{Third paragraph: reasons for these problems}
\begin{itemize}
    \item Reason 1: Heterogeneous preprocessing pipelines
    \item Reason 2: Inconsistent feature representations (voxel-wise vs ROI-based), not analysis cross microstructure features
    \item Reason 3: Evaluation protocols not directly comparable + results on different datasets
    \item Reason 4: Missing code
    \item Our solution: "To address these challenges, we propose: 1) an open-source and standardized benchmark; 2) of various machine learning techniques in dMRI; 3) on several tasks (classification, regression) 4) encompassing N datasets 
\end{itemize}

\noindent \textbf{Fourth paragraph: your contributions}
\begin{itemize}
    \item Standardized preprocessing pipeline for diffusion MRI
    \item Implemented various representation of the microstructural input data
    \item Unified feature extraction strategy
    \item Controlled evaluation protocol (train/test splits)
    \item Comparison of multiple feature representations
    \item Integration of both classical and DL predictors
    \item Public release of code and benchmark configuration
    \item End-to-end analysis
\end{itemize}

\section{Benchmark (Reuben)}
Start with a general introduction of this section that introduce each component of your benchmark + add a figure that presents an overview.

\subsection{Datasets}
\begin{itemize}
    \item Introduction sentence about why you picked these datasets
    \item Introduce Number of subjects
    \item Age range / demographics
    \item Scanner details
    \item Acquisition protocol or summary table with acquisition parameters
    \item Prediction targets (classification/regression)
\end{itemize}
If space (unlikely):
* PLOTS: Demographic distribution (age histograms) | Table with demographics summary (gender classes, age range...)

\subsection{Unified diffusion data preprocessing}
\subsubsection{Data preprocessing}
Following...., we propose to
\begin{itemize}
    \item Denoising
    \item Motion and eddy correction
    \item Spatial normalization
    \item Masking strategy
\end{itemize}

\subsubsection{Microstructural Representation}
Explain here the different forms of input representation.
The word "feature representation" is ambigiuous because neural network also extract features in a data-driven way. That's why I would rather use the term like "Microstructural Representation".
Explain why differences between white and gray matter. Give references here about why there are differences.

\paragraph{Gray Matter Representation}
\begin{itemize}
    \item Voxel-wise features
    \item High-dimensional representation (~64k features)
    \item Mask-based extraction
\end{itemize}
Explain that voxel-wise features preserve spatial resolution but increase dimensionality and risk of overfitting.

\paragraph{White Matter Representation}
\begin{itemize}
    \item Tract-based aggregation
    \item Mean values per tract (~48 features)
    \item Reduced dimensionality
    \item Anatomical interpretability
\end{itemize}
Justify tract-averaged representation: Noise reduction, improved stability, lower dimensional feature space, anatomically meaningful structure.

\subsubsection{Microstructural normalization}
    \begin{itemize}
        \item Binary targets (encoded 0/1), Continuous targets (standardized during training)
        \item Standardization fitted on training set only
        \item Applied to validation/test sets 
        \item Inverse transform for regression outputs (for visualization)
    \end{itemize}

\subsection{Prediction models}
    \begin{itemize}
        \item Dummy baseline (e.g., majority class for classification, mean predictor for regression)
        \item Classical ML pipeline: Linear models (e.g., logistic regression, linear regression), Random Forests, Support Vector Machines (SVM) + their PCA versions
        \item Deep Learning models (e.g., dinov2, medicalnet)
        \item Unified hyperparameter search strategy (e.g., grid search or random search with cross-validation on the training set)
\end{itemize}

\subsection{Benchmarks tasks and evaluation metrics}
Task selection rationale: Clinical relevance, data availability, and established benchmarks in the literature.

\subsubsection{Classification} 
\begin{itemize}
    \item Task description: Binary diagnosis prediction (e.g., ASD vs. TD)
    \item Evaluation Metrics: Accuracy, AUC, F1-score
\end{itemize}

\subsubsection{Regression} 
\begin{itemize}
    \item Task description: Age prediction (brain age estimation)
    \item Evaluation Metrics: R², MAE
\end{itemize}

\subsection{Evaluation protocols}

\begin{itemize}
    \item Data Splitting Strategy (Fixed train/test splits)
    \item Cross-validation within training set. Stratification
    \item No information leakage
    \item Mean ± std across folds
    \item Statistical comparison tests (e.g., paired t-tests, Wilcoxon signed-rank tests) to assess significance of performance differences between models
\end{itemize}


\section{Results (Reuben)}
Introduction sentence

\subsubsection{Choice of prediction models}
Table: 
For all tasks, choose one brain structure and one microstructural representation (choose widely this one) and report results in a table + comment

\subsubsection{Choice of the brain structure}
For one regression task and one classification task, show the impact of the brain structure with one miscrostructural representation. You can discard models that underperformed in the first model section.

\subsubsection{Choice of the microstructural representation}
For one task, show the impact of the microstructural representation for each brain structure. You can discard models that underperformed in the first model section.

\subsubsection{Training size impact} I think we will not have space for this, but here it would be for one regression task and one classification task, one structure, one microstrucral representation, impact of the model, here try to have variability (PCA, non-PCA, deep) for example.

\subsubsection{Methodological insights}
\begin{itemize}
    \item DL vs ML based
    \item Importance of dimensionality control
    \item Effect of tract-based aggregation 
    \item Risk of overfitting in mesh representations without PCA bcs of the high number of vertices representations
\end{itemize}

\section{Conclusion (Reuben)}
\begin{itemize}
    \item Introduced standardized dMRI benchmark for age prediction in healthy individuals using three large datasets (HCP, CamCAN, ABIDE) 
    \item Compared feature representations (MD, SH, MK, RTOP) and models (dummy classifier, linear, pca\_linear, forest, pca\_forest, svm, pca\_svm)
    \item Provided reproducible end{-}to{-}end framework and enables fair model comparison. Reduces methodological heterogeneity
    \item A step toward methodological unification/code availability
\end{itemize}


\newpage